\documentclass[10pt,a4paper]{amsart}
\usepackage[margin=19mm]{geometry}
\usepackage{amsmath,amssymb,mathtools}
\usepackage[scr=rsfs,cal=euler]{mathalfa}
\usepackage{xfrac}
\usepackage[unicode,hidelinks,bookmarks=false]{hyperref}
\newcommand{\ie}{i.e.\ }
\newcommand{\cf}{cf.\ }
\newcommand{\resp}{respectively\ }
\newcommand{\Subsection}{\subsection}
\providecommand{\nobreakdash}{\nobreak\hbox{-}}
\setlength{\emergencystretch}{2em}
\multlinegap0pt
\usepackage{amssymb}
\usepackage{enumerate}
\newtheorem{lemma}{Lemma}
\title{A link between error terms when counting smooth and rough numbers}
\author{Andreas Weingartner}
\date{Source: arXiv:2604.22058v1 (CC BY 4.0)}
\begin{document}
\maketitle

\noindent\emph{Source and licence.} A link between error terms when counting smooth and rough numbers. Version arXiv:2604.22058v1 (CC BY 4.0), distributed by arXiv under the Creative Commons Attribution 4.0 International licence, http://creativecommons.org/licenses/by/4.0/, as recorded in the arXiv licence metadata for this exact version and shown on the abstract page https://arxiv.org/abs/2604.22058v1. Full licence notice: \texttt{licenses/arxiv-2604.22058-CC-BY-4.0.txt}.\par\smallskip

\noindent\textit{Editorial scope.} Counted unit: the article's lemma LemRQ (source0.tex lines 318--335). The article's complete original proof of it is delivered with it (source0.tex lines 336--346, one \texttt{proof} environment of 11 lines). No mathematics is written, repaired or re-derived: the statement and the proof are the article's own lines, in the article's own notation, and no retained line is substituted or re-flowed.  Retained uncounted as the source-local context the proof needs: lines 274--285; lines 291--296; lines 307--312.  Nothing was omitted from any retained span, so the package carries no omission record.  No retained line contains a citation, so the wrapper carries no bibliography and no bibliography entry.  No retained line contains a figure environment, an image-inclusion command or drawing code, so no image is delivered and the package declares no asset dependency.  Editorial formatting only, in addition: an AMS \texttt{amsart} preamble that restates the article's own declarations and macros used by the retained lines, namely the environments \texttt{lemma} on separate counters, together with the standard packages those lines need.  The retained environments print in this document as Lemma 1 in order of appearance, whereas the article prints the same items with the numbering of its own full document.  The complete native source of the article as distributed in the arXiv e-print for this exact version is bundled unchanged as \texttt{sources/199067/source0.tex}.
\begin{lemma}\label{Lem0}
For $0< x \le y$, $y\ge 2$, we have
$$
Q^*(x,y)=0, \quad R^*(x,y)=0,
$$
$$
Q(x,y) = x\left( \Pi(y) - \frac{e^{-\gamma}}{\log y} \right), 
$$
$$
R(x,y) =  x\left( 1 - \frac{e^{-\gamma}}{\log y}\Pi(y)^{-1} \right) .
$$
\end{lemma}

\begin{lemma}\label{LemMobInv}
For $x> 0$, $y\ge 2$, we have
\begin{equation*}
Q(x,y)= \sum_{n\in S_y} \mu(n) R(x/n,y), \quad Q^*(x,y)= \sum_{n\in S_y} \mu(n) R^*(x/n,y).
\end{equation*}
\end{lemma}

\begin{lemma}\label{PsiLem}
Let $A>0$ and $\delta>0$ be fixed. We have
$$
\Psi(x,y) \ll x e^{-Au} \qquad (x\ge 1, \ y\ge e^{A+\delta}).
$$
\end{lemma}

\begin{lemma}\label{LemRQ}
Let $A>0$ and $\delta>0$ be fixed. If
$$
Q(x,y) \ll x f(y) e^{-A u} \qquad (x\ge 1,\  y\ge e^{A+\delta})
$$
then
$$
R(x,y)  \ll x f(y) e^{-A u}  \log y \qquad (x\ge 1,\  y\ge e^{A+\delta}).
$$
 If
$$
R(x,y) \ll x f(y) e^{-A u} \qquad (x\ge 1,\  y\ge e^{A+\delta})
$$
then
$$
Q(x,y)  \ll x f(y) e^{-A u}  \log y \qquad (x\ge 1,\  y\ge e^{A+\delta}).
$$
\end{lemma}

\begin{proof}
Assume the first upper bound for $Q(x,y)$. When $0<x<1$, this bound still holds, 
by Lemma \ref{Lem0} and since it holds when $x=1$. Thus 
$$
\frac{R(x,y)}{x f(y)} \ll \sum_{n\in \mathcal{S}_y} \frac{ \exp(-A\log(x/n)/\log y)}{n}
= e^{-Au} \sum_{n\in \mathcal{S}_y} \frac{\exp(A\log(n)/\log y)}{n }.
$$
Abel summation and Lemma \ref{PsiLem}, with $(A+\delta/2,\delta/2)$ in place of $(A,\delta)$, 
shows that the last sum is $\ll \log y$, which implies the desired estimate for $R(x,y)$. 
The second part is proved similarly, with the help of Lemma \ref{LemMobInv}.
\end{proof}

\end{document}
